\subsection{RW variance for chained conditional imputation models}
\label{extension-1}
The RW framework of Section~\ref{robinswang-variance-components} requires
three ingredients
tied to a joint parametric imputation likelihood: a parameter vector
\(\psi\) indexing the joint imputation distribution, an observed data
score equation \(\sum_i S_{\mathrm{obs},i}(\psi)=0\) whose solution
defines \(\hat\psi\), and the resulting missing value score
\(S_{\mathrm{mis},i}\) and nuisance influence \(d_i\).  MICE supplies
none of these.  Its conditional working models are fitted block by
block, where a block is a single variable imputed by its own
conditional model (the default in \texttt{mice}) or a group of
variables imputed jointly; the blockwise fits do not correspond to any
joint imputation likelihood, and admit no global parameter solving a
single score equation.  Even if one concatenated the blockwise
parameters as \(\psi = (\psi_1,\ldots,\psi_R)\), where \(R\) denotes
the number of imputation blocks, MICE conditional models are in
general mutually incompatible: they need not correspond to any joint
distribution over the missing variables, so no joint likelihood exists
and no global score equation of the form
\(\sum_i S_{\mathrm{obs},i}(\psi)=0\) can be defined.
Our proposed resolution is to shift from one global estimated parameter
to one realized parameter draw per block per imputation.  In each
imputation \(p\), MICE generates a parameter state \(\psi_r^{(p)}\) for
block \(r\), the coefficient draw used to produce that block's completed
values, and this realized state is already recorded in a standard MICE
run in \texttt{mice}.  We evaluate the RW score and nuisance components
at these realized states, one block at a time, and stack the per-block
contributions into composite vectors that are substituted into the
unchanged RW formula.  No joint imputation likelihood is assumed, and
no modeling choices are required beyond the conditional models that
MICE already fits.
The non-trivial step is justifying that stacking works.  Two conditions
must hold: the realized states \(\psi_r^{(p)}\) must concentrate around
their blockwise working limits \(\psi_{r0}\) fast enough that evaluating
at \(\psi_r^{(p)}\) rather than \(\psi_{r0}\) changes the empirical
moments only at negligible order; and the stacked empirical moments must
consistently estimate the population RW variance components for the MICE
estimator without a coherent joint likelihood.  Both are established in
Appendix~\ref{appendix-theory}; the rest of this section builds the
construction.
\subsubsection{Sequential conditional working law}
Suppose the missing variables are imputed in blocks \(r=1,\ldots,R\).
For block \(r\), let \(H_{r,i}^{(p)}\) denote the predictors and
previously updated variables used by the conditional model in completed
dataset \(p\), with working model
\(f_r\{Y_{r,i}\mid H_{r,i}^{(p)};\psi_r\}\).
Because the conditional models need not be mutually compatible, each
block has its own working parameter \(\psi_r\), whose working value
\(\psi_{r0}\) is defined as the probability limit of the corresponding
conditional estimator, rather than as a coordinate of a global
\(\psi\).  There is one \(\psi_r\) per block and no relationship between
the \(\psi_r\) values across blocks is assumed or needed.
In a proper MICE run, imputation \(p\) draws a parameter state
\(\psi_r^{(p)}\) for each block and generates the completed block from
it.  These realized states are already stored as part of the imputation
output in \texttt{mice}.
We evaluate all RW score components at these realized states
\(\psi_r^{(p)}\).  They are the natural plug-in: already produced by the
MICE run, requiring no additional estimation, and under standard
regularity conditions for proper imputation they concentrate around the
blockwise working limits at rate \(n^{-1/2}\),
\[
\psi_r^{(p)}-\psi_{r0}=O_p(n^{-1/2}).
\]
Because the RW score components are smooth functions of the imputation
parameter, a first-order Taylor expansion shows that evaluating at
\(\psi_r^{(p)}\) rather than \(\psi_{r0}\) changes the empirical RW
moments only at order \(n^{-1/2}\), which is higher order than the
leading term that determines the sandwich limit.  The substitution error
is therefore asymptotically negligible.  The formal statement is in
Appendix~\ref{appendix-theory}.
Next, we describe the score and nuisance objects required for the RW
variance computation.  Write the conditional-model score contribution
for block \(r\), subject \(i\), imputation \(p\) as
\[
s_{r,i}^{(p)}(\psi_r)
=
\frac{\partial}{\partial\psi_r}
\log f_r\{Y_{r,i}\mid H_{r,i}^{(p)};\psi_r\}.
\]
The missing-value score for rows imputed in block \(r\) is the working
score evaluated at the realized state,
\(S_{r,\mathrm{mis},i}^{(p)}=s_{r,i}^{(p)}(\psi_r^{(p)})\); for rows
not imputed in block \(r\), it is zero.  A large score means a small
change in \(\psi_r\) would have moved the imputed value substantially.
The nuisance influence is constructed as follows.  Let \(\delta_{r,i}=1\)
indicate that block \(r\) is observed for subject \(i\), and set
\(S_{r,\mathrm{obs},i}^{(p)}=\delta_{r,i}\,s_{r,i}^{(p)}(\psi_r^{(p)})\).
Define
\[
A_r^{(p)}
=
P_n\!\left\{
\left.
\frac{\partial s_{r,i}^{(p)}(\psi_r)}{\partial\psi_r^\top}
\right|_{\psi_r=\psi_r^{(p)}}
\delta_{r,i}
\right\},
\]
the observed-data Hessian of the block-\(r\) conditional log-likelihood.
The blockwise nuisance influence is the Newton linearization
\[
d_{r,i}^{(p)}
=
-\{A_r^{(p)}\}^{-1}S_{r,\mathrm{obs},i}^{(p)},
\]
the block-\(r\) analogue of the single-model \(d_i\): it measures how
much the block-\(r\) parameter draw would shift if subject \(i\)'s data
changed.  It requires no additional modeling beyond the conditional model
that MICE already fits.  Like the single-model case, \(d_{r,i}^{(p)}\)
is supported only on subjects with block \(r\) observed.
Stacking across all \(R\) blocks then gives the MICE imputation score and
nuisance influence vectors for subject \(i\) in imputation \(p\):
\[
S_{\mathrm{mis},i}^{(p)}
=
\bigl(
S_{1,\mathrm{mis},i}^{(p)\top},\ldots,S_{R,\mathrm{mis},i}^{(p)\top}
\bigr)^\top,
\qquad
d_i^{(p)}
=
\bigl(
d_{1,i}^{(p)\top},\ldots,d_{R,i}^{(p)\top}
\bigr)^\top.
\]
The conditions required of each block are minimal: a regular
observed-score derivative (so that \(A_r^{(p)}\) is invertible), finite
second moments for the score contributions, and a realized state that
concentrates at its working limit.  No coherent joint likelihood across
blocks is needed.
Stacking automatically captures cross-block parameter uncertainty.  The
off-diagonal elements of \(\alpha = E[d_i d_i^\top]\), such as
\(E[d_{r,i}\,d_{r',i}^\top]\) for \(r\neq r'\), reflect correlation
between estimation uncertainty in different blocks, which arises because
the same observed data contribute to both conditional model fits.  These
terms are non-zero in general and are automatically included when
\(\widehat\alpha\) is formed from the full stacked vectors \(d_i^{(p)}\).
\subsubsection{RW components for the MICE estimator}
The analysis estimator is still defined by the averaged completed-data
estimating equation \(0=P_n\bar U(\hat\beta)\).  All RW components are
therefore similar to their Section~\ref{robinswang-variance-components}
counterparts, with the single-model \(S_{\mathrm{mis},i}^{(p)}\) and
\(d_i^{(p)}\) replaced by the stacked MICE quantities defined above; we
call the resulting variance estimator RW-MICE.
Concretely, the bread \(\hat\tau\) is computed exactly as before; the
three covariance pieces are:
\begin{align*}
\widehat\Omega
&=
\frac{1}{n}\sum_{i=1}^n \bar U_i(\hat\beta)\,\bar U_i(\hat\beta)^\top,
\\[4pt]
\widehat\kappa
&=
\frac{1}{nm}\sum_{p=1}^m\sum_{i=1}^n
U_i^{(p)}(\hat\beta)\,S_{\mathrm{mis},i}^{(p)\top},
\qquad
\widehat\alpha
=
\frac{1}{nm}\sum_{p=1}^m\sum_{i=1}^n
d_i^{(p)}d_i^{(p)\top},
\\[4pt]
\widehat C
&=
\frac{1}{n}
\bigl[
\widehat\kappa\,\bar d^\top \bar U
+
\bar U^\top\bar d\,\widehat\kappa^\top
\bigr],
\end{align*}
where \(\bar U\) and \(\bar d\) are the \(n\times q\) and
\(n\times\dim(\psi)\) matrices with rows \(\bar U_i^\top\) and
\(\bar d_i^\top\), and \(\dim(\psi)=\sum_r\dim(\psi_r)\) is the total
number of imputation parameters across all blocks.
Stacking changes the internal block structure of the matrices.
The matrix \(\widehat\kappa = [\widehat\kappa_1,\ldots,\widehat\kappa_R]\)
is block-partitioned by imputation block, with
\(\widehat\kappa_r = (nm)^{-1}\sum_{p,i}U_i^{(p)}S_{r,\mathrm{mis},i}^{(p)\top}\)
measuring how sensitive the analysis score is to block \(r\)'s parameter;
because the analysis equation mixes contributions from all imputed
variables, \(\widehat\kappa_r\) is typically non-zero even for blocks
whose variables do not appear directly in the analysis estimating
function (in parametric regression analysis models, blocks other than
the primary analysis covariate).  The matrix
\(\widehat\alpha\) has \((r,r')\) block
\(\widehat\alpha_{rr'} = (nm)^{-1}\sum_{p,i}d_{r,i}^{(p)}\,d_{r',i}^{(p)\top}\);
the off-diagonal blocks for \(r\neq r'\) capture cross-block parameter
uncertainty and are non-zero whenever subjects contribute to both
conditional model fits.  Dropping them would incorrectly assume
independence of estimation uncertainty across blocks, which fails unless
missing data patterns are block-monotone with no shared covariates.
Consistency of the stacked estimator is established in
Appendix~\ref{appendix-theory}; when \(R=1\) the construction reduces
to the single-model RW estimator of
Section~\ref{robinswang-variance-components}.
The MICE extension requires no new modeling beyond what a standard MICE
run already produces.  Given the completed datasets and the recorded
parameter draws \(\psi_r^{(p)}\), every RW component is computed from
formulas unchanged in form from Section~\ref{robinswang-variance-components}.
The only structural change is that \(S_{\mathrm{mis}}\) and \(d\) are
block vectors, and the resulting sandwich automatically accounts for
within-block and cross-block parameter uncertainty through the full
block structure of \(\widehat\alpha\).